\section{Simulation in the phase-polynomial representation}
\label{sec:simulator-phase-polynomial-representation}

We now describe how to update the phase-polynomial representation under the
operations in a 3PP circuit.  We represent the current state as
\begin{equation}
    \ket{\psi}
    = \frac{1}{\sqrt{|S|}}
      \sum_{x\in S}\omega_8^{q(x)}\ket{x},
    \qquad S=x_0\oplus V,
    \label{eq:phase-polynomial-state-update}
\end{equation}
where $S\subseteq\F_2^n$ is affine and $q$ is a third-order phase
polynomial.  For the initial state
$\ket{0}^{\otimes a}\ket{+}^{\otimes b}$, one takes
$S=\{0\}^a\times\F_2^b$ and $q=0$.
Thus, the representation can be initialized and all the updates below can be carried out in polynomial time by linear algebra over $\F_2$ and arithmetic on the coefficients of $q$.

\subsection{Unitary operations}

Every CNOT-dihedral unitary has an affine permutation and a third-order
phase polynomial $p$ such that
\begin{equation}
    U\ket{x}=\omega_8^{p(x)}\ket{\phi(x)},
    \qquad \phi(x)=Ax\oplus b,
    \label{eq:cnot-dihedral-action}
\end{equation}
where $A\in\F_2^{n\times n}$ is invertible.  Applying $U$ to
\cref{eq:phase-polynomial-state-update} gives the updated data
\begin{equation}
    S'=\phi(S),
    \qquad
    q'(y)=q\bigl(\phi^{-1}(y)\bigr)
           +p\bigl(\phi^{-1}(y)\bigr)\pmod 8.
    \label{eq:phase-polynomial-unitary-update}
\end{equation}
The set $S'$ is affine, and $q'$ is again a third-order phase polynomial by
\Cref{cor:phase-polynomial-affine-pullback}.  In particular, the elementary
updates are as follows, where
$\tau_{a\to b}(x):=x\oplus x_a e_b$:
\begin{align}
    \T_a^{\pm1}:&\quad
        S'=S,
        &q'(y)&=q(y)\pm y_a, \\
    Z_a:&\quad
        S'=S,
        &q'(y)&=q(y)+4y_a, \\
    X_a:&\quad
        S'=S\oplus e_a,
        &q'(y)&=q(y\oplus e_a), \\
    \CNOT_{a\to b}:&\quad
        S'=\tau_{a\to b}(S),
        &q'(y)&=q\bigl(\tau_{a\to b}(y)\bigr).
\end{align}
Updates for arbitrary Pauli operators follow by composing the $X_a$ and
$Z_a$ rules. Global phases may be discarded.

\subsection{Z-measurements}
\label{sec:z-measurements-3pp}

Let us consider a measurement of a $Z$-type Pauli operator $Z(v)=\prod_{a=1}^n Z_a^{v_a}$, and encode its measurement outcome
as $(-1)^\mu$, with $\mu\in\F_2$.
The projection onto that outcome can be written
\begin{align}
    \ket{\psi'_\mu} &= \frac{1}{2}\left( 1 + (-1)^\mu Z(v) \right)\ket{\psi} \\
    &= \frac{1}{2\sqrt{|S|}} \sum_{x\in S}
    \left( 1 + (-1)^{\mu+\langle v,x\rangle} \right)\ket{x} \\
\end{align}
where we used $Z(v)\ket{x}=(-1)^{\langle v,x\rangle}\ket{x}$.
Therefore, the phase is only nonzero when $\langle v,x\rangle=\mu$, and the affine support is restricted to that hyperplane:
\begin{equation} \label{eq:phase-polynomial-z-measurement-update}
    S'_\mu
    =S\cap\{x\in\F_2^n:\langle v,x\rangle=\mu\},
    \qquad
    q'_\mu=q\big|_{S'_\mu}.
\end{equation}
The probability of the outcome is
\begin{equation}
    \Pr(\mu)=\frac{|S'_\mu|}{|S|}.
\end{equation}
The linear function $x\mapsto\langle v,x\rangle$ is either constant on
$S$ or balanced on $S$.  Consequently, the two outcome probabilities are
either $0$ and $1$, or both $1/2$.  Gaussian elimination determines which
case applies and produces an affine description of $S'_\mu$ in polynomial
time.  A single-qubit $Z$-measurement is the special case $v=e_a$.

\subsection{X-measurements}

Recall from \Cref{sec:compatible-measurements} that an $X$-type Pauli measurement is compatible with a 3PP state when every outcome of nonzero probability leaves a 3PP state. We first introduce the discrete derivative and classify single-qubit $X$ measurements with their phase-polynomial updates, then extend the results to arbitrary $X$-type Pauli measurements.

\begin{definition}[Discrete derivative]\label{def:discrete-derivative}
    Given a function $f: V \to \mathbb Z_8$, we define its discrete derivative $\Delta_v f$ w.r.t. $v \in V$ by $\Delta_v f(x) := f(x \oplus v) - f(x) \pmod 8.$
\end{definition}


\newcommand{\xMeasurementClassificationStatement}{%
    Consider a 3PP state
    \[ \ket{\psi} = \frac{1}{\sqrt{|S|}} \sum_{x\in S} \omega_8^{q(x)}\ket{x} \]
    over an affine subspace $S = x_0 \oplus V$ with phase polynomial $q$.
    An $X$ measurement on qubit $a$ is compatible with $\ket{\psi}$ iff one of the following cases apply.
    In each case, we get a post-measurement affine subspace $S'_\mu \subseteq \mathbb F_2^n$ and phase polynomial $q'_\mu: S'_\mu \to \mathbb Z_8$ depending on the measurement outcome $\mu \in \{0, 1\}$, which occurs with some probability $Pr(\mu)$.
    \begin{enumerate}
        \item $e_a \not\in V$. In that case, $Pr(\mu) = \frac{1}{2}$ and the post-measurement state is defined by $S'_\mu = S \cup (S \oplus e_a)$ and the unique phase polynomial $q'_\mu$ with $q'_\mu(x) = q(x)$ and $q'_\mu(x \oplus e_a) = q(x) + 4\mu$ for all $x \in S$.

        \item
        $e_a \in V$ and one of the following apply:
        \begin{enumerate}
            \item \label{thm:x-measurements:even-0-4}
            $\Delta_{e_a} q(x) = 4f(x) \pmod 8$ for some affine function $f: S \to \mathbb F_2$. In that case
            \[
                Pr(\mu) = \begin{cases}
                    0 & \text{if } f(x) \neq \mu \text{ for all } x \in S, \\
                    1 & \text{if } f(x) = \mu \text{ for all } x \in S, \\
                    \frac{1}{2} & \text{otherwise},
                \end{cases}
            \]
            and the post-measurement state is defined by $S'_\mu = S \cap \{ x ~|~ f(x) = \mu \pmod 2 \}$ and $q'_\mu(x) = q(x)$ for all $x \in S'_\mu$.

            \item \label{thm:x-measurements:even-2-6}
            $\Delta_{e_a} q(x) = 2 + 4g(x) \pmod 8$ for some function $g: S \to \mathbb F_2$ such that $2g$ is a third-order phase polynomial.
            In that case, $Pr(\mu) = \frac{1}{2}$ and the post-measurement state is defined by $S'_\mu = S$ and $q'(x) = q(x) - 2(-1)^\mu g(x)$

            \item \label{thm:x-measurements:odd}
            $\Delta_{e_a} q(x) = (-1)^{h(x)}d \pmod 8$ for some affine function $h: S \to \mathbb F_2$ and $d \in \{ 1, 3, 5, 7 \}$.
            In that case, $Pr(\mu) = \frac{1}{2}(1 + (-1)^\mu\cos(d\pi/4))$ and the post-measurement state is defined by $S'_\mu = S$ and $q'(x) = q(x) + (4\mu - d)h(x)$.
        \end{enumerate}
    \end{enumerate}
}
\begin{theorem}\label{thm:x-measurements}
    \xMeasurementClassificationStatement
\end{theorem}
The proof can be found in \cref{appendix:x-measurements}.

For a general $X$-type Pauli operator $X(m)=\prod_j X_j^{m_j}$ with $m\in\F_2^n\setminus\{0\}$, choose $a$ with $m_a=1$ and set $C_m=\prod_{j\ne a:\,m_j=1}\CNOT_{a\to j}$. Then $C_mX(m)C_m^\dagger=X_a$, so
\begin{align}
    \frac{I+(-1)^\mu X(m)}{2}
    =C_m^\dagger\frac{I+(-1)^\mu X_a}{2}C_m.
    \label{eq:x-measurement-reduction}
\end{align}
Since both $C_m$ and $C_m^\dagger$ preserve 3PP states, measuring $X(m)$ on $\ket{\psi}$ has the same compatibility and outcome probabilities as measuring $X_a$ on $C_m\ket{\psi}$. Applying $C_m^\dagger$ to the conditional state gives the original measurement's conditional state. In the original coordinates, the same cases and update formulas apply with $e_a$ replaced by $m$ and $\Delta_{e_a}q$ replaced by $\Delta_mq$; in particular, case~(1) means $m\notin V$. We use these case labels for arbitrary $X$-type Pauli measurements throughout the paper. For $m=0$, the operator is the identity and the outcome is deterministically $+1$.

\subsection{Efficient simulation}

% Supply the representation's equation label only in the main-text statement.
\newcommand{\efficientSimulationStatement}[1]{%
An $n$-qubit 3PP state admits a representation
\begin{align}
    \ket{\psi}
    = 2^{-r/2}\sum_{y\in\mathbb{F}_2^r}
    \omega_8^{p(y)}\ket{x_0\oplus By},
    #1
\end{align}
where $B\in\mathbb{F}_2^{n\times r}$ has full column rank and $p:\mathbb{F}_2^r\to\mathbb{Z}_8$ is a third-order phase polynomial.
This representation requires $O(n^3)$ bits and supports the following operations.

\begin{enumerate}
    \item \textbf{Circuit simulation.}
    Let $C$ be a 3PP circuit made of $m$ operations acting on at most $n$ qubits.
    There exists a classical algorithm that samples the measurement record according to its exact distribution and maintains the corresponding conditional state, up to global phase, in the form \eqref{eq:simulation-representation}.
    Each operation takes $O(n^3)$ expected time, giving a total expected runtime of
    \begin{align}
        O(mn^3)+T_{\mathrm{classical}},
    \end{align}
    where $T_{\mathrm{classical}}$ is the total time required for classical feedforward.
    The maintained state representation uses $O(n^3)$ bits, excluding storage for the measurement record and classical processing.

    \item \textbf{Queries on represented states.}
    Given a state in the form \eqref{eq:simulation-representation}, the following tasks can be performed:
    \begin{enumerate}
        \item Compute $\langle x|\psi \rangle$ exactly for any $x\in\mathbb{F}_2^n$ in $O(n^3)$ time.
        \item Sample $x\in\mathbb{F}_2^n$ according to $|\langle x|\psi \rangle|^2$ in $O(n^2)$ time.
        \item Compute $\bra{\psi}P\ket{\psi}$ exactly for any specified Pauli operator $P$ in $O(n^3)$ time.
        \item Given a second represented state on the same ordered qubits, decide whether the two states are equal up to global phase in $O(n^4)$ time.
    \end{enumerate}
\end{enumerate}
}

\begin{theorem}[Efficient simulation of 3PP circuits]
\label{thm:efficient-simulation}
\efficientSimulationStatement{\label{eq:simulation-representation}}
\end{theorem}
The proof can be found in \cref{appendix:efficient-simulation}.

The Pauli-expectation claim does not require the
corresponding measurement to be compatible: evaluating
an expectation value does not require representing the
post-measurement state. By linearity, it also gives
efficient exact expectations for observables supplied
as polynomial-size sums of Pauli operators.

Equality testing should be distinguished from evaluating
an arbitrary overlap between two DCP states. Such an
overlap generally involves a third-order phase sum,
rather than the quadratic sum appearing in
\cref{eq:pauli-expectation-gauss-sum}, and is not
covered by the theorem.

For circuits with stochastic Pauli noise, the theorem
applies to each sampled fault trajectory, provided that
all measurements remain compatible along that trajectory.
The represented state and its exact expectation values
are conditional on the sampled faults and measurement
outcomes.

Our simulation theorem recovers efficient simulation of Clifford circuits. Every Pauli stabilizer state is a 3PP state, and Pauli measurements preserve the stabilizer-state class. While $\CNOT$ and $\S=\T^2$ are already supported, a Hadamard gate can be implemented by the measurement-based gadget shown in \cref{fig:hadamard-gadget}.

For stabilizer inputs, every intermediate state remains a stabilizer state, so this measurement is compatible. Thus any Clifford circuit can be compiled into a 3PP circuit with constant overhead per gate.

Finally, we note that alternative sparse representations of the phase polynomial (for example in the form of Eq. \eqref{eqn:phase-polynomial-def}), can yield better complexity for circuits with few non-Clifford gates, which we have exploited in our implementation.
